\Needspace{6\baselineskip}
\begin{center}\hypertarget{native-252-article180}{}\pdfbookmark[2]{180.}{article-180}180.\end{center}
The principle set out in the preceding article also has the advantage that the numerical determination of the unknowns reduces to a very straightforward algorithm whenever the functions $V$, $V'$, $V''$ etc. are linear. Suppose that
\[
\begin{gathered}
M-V=v=-m+ap+bq+cr+ds+\text{etc.}\\
M'-V'=v'=-m'+a'p+b'q+c'r+d's+\text{etc.}\\
M''-V''=v''=-m''+a''p+b''q+c''r+d''s+\text{etc.}\\
\text{etc.,}
\end{gathered}
\]
and let us set
\[
\begin{gathered}
av+a'v'+a''v''+\text{etc.}=P\\
bv+b'v'+b''v''+\text{etc.}=Q\\
cv+c'v'+c''v''+\text{etc.}=R\\
dv+d'v'+d''v''+\text{etc.}=S\\
\text{etc.}
\end{gathered}
\]
Then the $\nu$ equations of Article 177 from which the values of the unknowns are to be determined are clearly
\[
P=0,\qquad Q=0,\qquad R=0,\qquad S=0\quad\text{etc.,}
\]
provided we suppose the observations to be equally good, the case to which we showed in the preceding article how to reduce the others. There are thus as many linear equations as there are unknowns to be determined, and their values can be found by ordinary elimination.

Let us now see whether this elimination is always possible, or whether the \SourceLeaf{253}solution can sometimes be indeterminate or even impossible. From the theory of elimination we know that the second or third case occurs when, after omitting one of the equations $P=0$, $Q=0$, $R=0$, $S=0$ etc., we can form from the others an equation either identical to the omitted one or inconsistent with it. Equivalently, this happens when we can assign a linear function
\[
\alpha P+\beta Q+\gamma R+\delta S+\text{etc.,}
\]
which is identically $=0$ or at least independent of all the unknowns $p$, $q$, $r$, $s$ etc. Suppose, then, that $\alpha P+\beta Q+\gamma R+\delta S+\text{etc.}=\varkappa$. We have the identity
\[
\begin{aligned}
&(v+m)v+(v'+m')v'+(v''+m'')v''+\text{etc.}\\
&\qquad=pP+qQ+rR+sS+\text{etc.}
\end{aligned}
\]
If we now suppose that the substitutions $p=\alpha x$, $q=\beta x$, $r=\gamma x$, $s=\delta x$ etc. turn the functions $v$, $v'$, $v''$ etc. respectively into $-m+\lambda x$, $-m'+\lambda'x$, $-m''+\lambda''x$ etc., we clearly have the identity
\[
\begin{aligned}
&(\lambda\lambda+\lambda'\lambda'+\lambda''\lambda''+\text{etc.})xx\\
&\qquad-(\lambda m+\lambda'm'+\lambda''m''+\text{etc.})x=\varkappa x,
\end{aligned}
\]
that is, $\lambda\lambda+\lambda'\lambda'+\lambda''\lambda''+\text{etc.}=0$, $\varkappa+\lambda m+\lambda'm'+\lambda''m''+\text{etc.}=0$. It follows necessarily that $\lambda=0$, $\lambda'=0$, $\lambda''=0$ etc. and $\varkappa=0$. Hence all the functions $V$, $V'$, $V''$ etc. are such that their values stay the same if the quantities $p$, $q$, $r$, $s$ etc. are given any increments or decrements proportional to the numbers $\alpha$, $\beta$, $\gamma$, $\delta$ etc. We have already observed above that cases of this kind fall outside our inquiry, since it is clear that the unknowns could not be determined even if the true values of the functions $V$, $V'$, $V''$ etc. were given.

We can readily reduce all remaining cases in which the functions $V$, $V'$, $V''$ etc. are nonlinear to the case considered here. Let $\pi$, $\chi$, $\rho$, $\sigma$ etc. denote approximate values of the unknowns $p$, $q$, $r$, $s$ etc. (which we readily obtain if from the $\mu$ equations $V=M$, $V'=M'$, $V''=M''$ etc. we at first use only $\nu$). In place of the unknowns we introduce new ones $p'$, $q'$, $r'$, $s'$ etc., putting $p=\pi+p'$, $q=\chi+q'$, $r=\rho+r'$, $s=\sigma+s'$ etc. The values of these new unknowns will clearly be so small that we may neglect their squares and products, so that the equations at once become linear. If, contrary to expectation, the completed calculation yields values of the unknowns $p'$, $q'$, $r'$, $s'$ etc. \SourceLeaf{254}so large that it seems unsafe to have neglected the squares and products, the difficulty can readily be remedied by repeating the same operation, taking, in place of $\pi$, $\chi$, $\rho$, $\sigma$ etc., the corrected values of $p$, $q$, $r$, $s$ etc.
